\chapter{Esecuzione, verifica e consegna}
\label{chap:esecuzione}

\section{Mappa della configurazione}

Prima dell'esecuzione è utile distinguere le variabili di ambiente dai
parametri globali del progetto. Il file \filepath{.env} locale contiene i
valori specifici della macchina e della lavorazione corrente; non deve
necessariamente ripetere tutte le opzioni facoltative presenti in
\filepath{.env.template}. Quando una di queste opzioni è assente, viene usato
il valore predefinito dichiarato in \filepath{graph_reconstruction/config.py}.

\begin{contracttable}
\envvar{IMAGES_DIR}, \envvar{LABELS_EXT_DIR},
\envvar{YOLO_MODEL_PATH} & Variabili di ambiente in \filepath{.env} &
Percorsi degli input, delle label estese e del checkpoint YOLO. Sono richiesti
dal caricamento della configurazione. \\

\envvar{WORK_DIR}, \envvar{ASSETS_DIR},
\envvar{ANNOTATIONS_DIR} & Variabili di ambiente in \filepath{.env} &
Directory degli output, delle risorse grafiche e delle annotazioni di grafo. \\

\envvar{DEBUG_OUTPUT_ENABLED}, \envvar{DEBUG_OUTPUT_DIR} & Variabili di
ambiente in \filepath{.env} & Abilitazione e destinazione dei file diagnostici
facoltativi. \\
\midrule

Opzioni VGGT elencate in \filepath{.env.template} & Variabili di ambiente
facoltative & Se non sono riportate nel \filepath{.env} locale, la pipeline usa
i valori predefiniti di \filepath{config.py}, compreso il percorso
dell'interprete VGGT. \\
\midrule

\code{CONF_THRESHOLD}, \code{YOLO_CONF}, \code{YOLO_IOU},
\code{YOLO_IMGSZ} & Parametri globali in \filepath{config.py} & Soglie e
dimensione di inferenza; non vengono lette dal file \filepath{.env}. \\

\code{SKIP_NODES}, \code{VGGT_SKIP_NODES}, \code{MIN_IMAGES} & Parametri
globali in \filepath{config.py} & Filtri applicati rispettivamente al batch,
all'elaborazione VGGT e al numero minimo di viste. \\
\midrule

Tabelle dei componenti e vincoli strutturali & Dati globali in
\filepath{config.py} & Comprendono codici, valenze, diametri, co-occorrenze,
triplette e coppie vietate; non sono impostazioni operative del singolo run. \\
\end{contracttable}

\Needspace{10\baselineskip}
\section{Elaborazione di un singolo nodo}

Per validare un nuovo dataset si consiglia di iniziare da un ID con viste
chiare. Per il nodo \code{7}:

\begin{lstlisting}[language={}]
conda activate graph_rec
python run.py 7
\end{lstlisting}

\filepath{run.py} verifica che l'ID esista tra le immagini scoperte e passa le
viste del nodo all'orchestratore. Se l'ID non viene trovato, stampa il percorso
di \envvar{IMAGES_DIR} che ha controllato: è un'indicazione utile per capire se
il problema è nel file \filepath{.env} o nei nomi delle immagini.

L'exit code non è sufficiente a certificare il run. \filepath{run.py} può
terminare con codice zero anche se l'ID non è stato trovato o se nessuna vista
ha prodotto detection valide. Cerca nel log l'intestazione \code{NODE 7}, il
blocco \code{FINAL GRAPH} e il percorso dell'output, poi verifica che
\filepath{fused_graph.json} sia stato appena scritto.

La prima elaborazione può richiedere più tempo per la generazione dei file
VGGT. Prima di interromperla, verificare l'attività di CPU, GPU o disco e la
creazione degli artefatti nella cartella del nodo.

\subsection{Interpretazione del log}

Il log articola la lavorazione in tre step principali. Nello
\emph{Step 1} ogni immagine viene aperta, analizzata da YOLO e trasformata in un
insieme di osservazioni valide. Una riga come
\code{node7_0002.jpg -> 5 nodes} significa che, dopo i filtri, quella vista ha
prodotto cinque componenti utilizzabili.

Prima dello \emph{Step 2} possono comparire i messaggi VGGT. Se gli artefatti
\code{.npz} mancano, l'interprete indicato da
\envvar{VGGT_EMBEDDED_PYTHON} li genera; se sono già presenti, li riusa. Lo
Step 2 è la ReID: raggruppa le osservazioni delle varie fotografie in
\emph{scene node}, cioè nelle istanze fisiche che la pipeline ritiene presenti
nella scena.

Lo \emph{Step 3} costruisce gli archi a partire dagli indizi geometrici, prova a
collegare le componenti rimaste separate e associa i diametri. Il riepilogo
\code{FINAL GRAPH} indica quanti nodi e archi sono stati ottenuti, quanti nodi
sono isolati e se il grafo finale è connesso.

In fondo viene stampato il percorso della cartella di output. È il riferimento
più affidabile quando \envvar{WORK_DIR} non coincide con
\filepath{./output}.

\section{Output prodotti}

Con la configurazione del Capitolo~\ref{chap:dati}, il nodo \code{7} viene
scritto in:

\begin{center}
\filepath{D:/GraphReconstructionData/work/impianto_A/multiview/node7/}
\end{center}

La struttura tipica, dopo ricostruzione e controllo visivo, è questa:

\begin{lstlisting}[language={}]
work/impianto_A/
  detection_cache/
    node7_0001.json
    node7_0002.json
  multiview/
    node7/
      fused_graph.json
      scene_nodes.json
      bird_view.png
      vggt_geometry.json
      vggt/
        generation_request.json
        generation_summary.json
        summary.json
        node7_0001.npz
        node7_0002.npz
      fused_graph_plot.png
      fused_graph_qgis_combined.geojson
      standard_similarity.json
  standard_comparison/
    similarity_table.csv
    similarity_table.json
    similarity_table.html
    best_matches.csv
    report_completo.html
\end{lstlisting}

Non tutti i file compaiono con il solo \filepath{run.py}.
\filepath{fused_graph_plot.png}, il GeoJSON e i risultati di confronto vengono
aggiunti dagli entrypoint di visualizzazione e valutazione. Inoltre VGGT può
essere disabilitato o escluso per alcuni nodi; in quel caso i relativi file non
sono un requisito di successo.

\subsection{Output principale: \filepath{fused_graph.json}}

Questo file è destinato alle applicazioni a valle. Riporta il
numero totale di immagini, di nodi e di archi e indica se il grafo è connesso.
Seguono due liste, \code{nodes} ed \code{edges}.

Ogni elemento di \code{nodes} rappresenta un componente fisico. Oltre a codice
e nome, conserva la confidenza, le immagini che lo supportano, il grado nel
grafo e la valenza massima prevista. \code{support_count} è particolarmente
utile nel controllo: un componente osservato in più viste è in genere meglio
supportato di un componente derivato da una sola detection.

Gli elementi di \code{edges} descrivono le connessioni. I campi
\code{vggt_score}, \code{vggt_support_count} ed \code{evidence_sources}
spiegano quale evidenza ha sostenuto l'arco. Se \code{forced} vale
\code{true}, il collegamento è stato aggiunto dal refinement per rendere il
grafo connesso. Non è necessariamente sbagliato, ma merita più attenzione nel
controllo umano, soprattutto quando \code{forced_reason} indica un fallback
strutturale.

I diametri compaiono come \code{diameter_mm}. I campi di confidenza e conflitto
permettono di distinguere un valore sostenuto in modo coerente da un ramo sul
quale sono state osservate indicazioni discordanti.

\subsection{Tracciabilità ReID: \filepath{scene_nodes.json}}

\filepath{scene_nodes.json} è il ponte tra fotografie e grafo. Per ogni istanza
fisica riporta le immagini di supporto e gli identificativi delle osservazioni
che sono state fuse. Se nel grafo finale compaiono due componenti che sembrano
lo stesso oggetto, questo è il primo file da aprire: permette di vedere se la
ReID ha lasciato separate detection provenienti da viste diverse.

È un artefatto di tracciabilità, normalmente escluso dalle interfacce
applicative.

\subsection{Geometria e render di controllo}

I file nella directory \filepath{vggt/} sono dati numerici intermedi costosi da
generare. \filepath{vggt_geometry.json} raccoglie invece la geometria aggregata,
gli indizi di connessione e le eventuali correzioni di codice riutilizzabili nei
run successivi.

\filepath{bird_view.png} viene creato durante la ricostruzione quando esiste
geometria sufficiente per almeno due scene node. La sua assenza non indica da
sola un fallimento, soprattutto se VGGT è disabilitato. Quando esiste, offre una
vista rapida della geometria ma non sostituisce il render del grafo. Per
ottenere quest'ultimo eseguire:

\begin{lstlisting}[language={}]
python visualize.py 7
\end{lstlisting}

Il comando legge \filepath{fused_graph.json}, salva
\filepath{fused_graph_plot.png}, crea un GeoJSON con layout di visualizzazione e
prova a generare il report con lo standard più simile.

Nel PNG verificare prima numero e tipo dei componenti, quindi gli archi. Se i
componenti sono già duplicati o mancanti, il problema precede la
costruzione del grafo ed è probabilmente nella detection o nella ReID.

\section{Verifica del risultato}

\subsection{Valutazione con ground truth}

Se sono disponibili label estese e grafo annotato, eseguire:

\begin{lstlisting}[language={}]
python evaluate.py 7
\end{lstlisting}

Il terminale mostra true positive, false positive, false negative, precision,
recall, F1 e purity delle istanze. Confronta inoltre il grafo con gli standard
di riferimento e salva la tabella completa sotto
\filepath{<WORK_DIR>/standard_comparison/}.

Le metriche su istanze e archi vengono stampate nel terminale ma non sono
salvate automaticamente in un report dedicato. Per una lavorazione tracciabile
è quindi opportuno acquisire anche il log completo del comando. I file persistiti
da questa fase riguardano soprattutto il confronto con gli standard.

Se mancano le label, la parte supervisionata viene saltata. Il confronto
strutturale con gli standard può comunque essere eseguito sul grafo ricostruito,
perché non dipende dalle bounding box ground truth.

I report aggregati sotto \filepath{standard_comparison/} descrivono gli ID
passati nell'esecuzione corrente. Un comando limitato a un solo nodo può quindi
sostituire un report creato in precedenza per l'intero batch. Genera i report di
consegna soltanto dopo aver stabilito l'insieme definitivo di nodi.

\section{Esecuzione batch e preparazione della consegna}

\subsection{Elaborazione batch}

Dopo la validazione del nodo campione, \filepath{run.py} senza argomenti
elabora gli ID rilevati che non appartengono a \code{SKIP_NODES} e che hanno
almeno il numero di viste definito da \code{MIN_IMAGES}; entrambi i valori si
trovano in \filepath{graph_reconstruction/config.py}:

\begin{lstlisting}[language={}]
python run.py
\end{lstlisting}

Con Conda si raccomanda di eseguire i tre entrypoint nello stesso terminale:

\begin{lstlisting}[language={}]
python run.py
python evaluate.py
python visualize.py --all
\end{lstlisting}

Per limitare il flusso a un sottoinsieme, indicare gli stessi ID nei tre
comandi:

\begin{lstlisting}[language={}]
python run.py 2 7 15_1
python evaluate.py 2 7 15_1
python visualize.py 2 7 15_1
\end{lstlisting}

\subsection{Export per QGIS}

Per convertire tutti i grafi disponibili:

\begin{lstlisting}[language={}]
python json_adapter_to_qgis.py
\end{lstlisting}

Per scegliere la geometria disponibile su un solo nodo:

\begin{lstlisting}[language={}]
python json_adapter_to_qgis.py `
  D:/GraphReconstructionData/work/impianto_A/multiview/node7/fused_graph.json `
  --geometry-source auto
\end{lstlisting}

Il GeoJSON mantiene nodi, archi e attributi, ma le coordinate sono locali alla
scena oppure sintetiche. Prima di sovrapporlo a cartografia reale occorre una
georeferenziazione esterna.

\filepath{evaluate.py} e \filepath{visualize.py} generano entrambi un GeoJSON
con geometria \code{fake} e possono sovrascrivere un export creato in
precedenza con \code{--geometry-source auto} o \code{vggt}. Se il GeoJSON fa
parte della consegna, esegui l'export con la geometria desiderata come ultimo
passaggio del flusso.

\subsection{Cache e riesecuzioni}

La cache delle detection si trova in
\filepath{<WORK_DIR>/detection_cache/}; i file VGGT restano nella cartella del
nodo. Questa scelta rende i run successivi più veloci, ma richiede disciplina:
la cache di detection non riconosce ogni possibile cambio dell'immagine, del
modello o del decoder.

Quando cambia il dataset o il modello YOLO, utilizzare una nuova
\envvar{WORK_DIR}. Questa scelta evita cancellazioni selettive e permette di
confrontare le due lavorazioni senza perdere traccia di quale cache sia stata
usata.

\subsection{Checklist prima del batch completo}

Prima di avviare un batch completo verifica che:

\begin{itemize}
  \item l'ambiente Conda corretto sia attivo e il decoder sia importabile;
  \item \filepath{.env} punti al dataset, ai pesi e alla directory di lavoro
        della commessa corrente;
  \item i nomi delle immagini rispettino il pattern e un singolo ID venga
        scoperto correttamente;
  \item il volume di \envvar{WORK_DIR} abbia spazio per detection, point map e
        report;
  \item il nodo campione produca \filepath{fused_graph.json} e un render
        leggibile;
  \item gli archi forzati del campione siano stati controllati;
  \item le metriche siano disponibili oppure sia registrato il motivo per cui
        il ground truth manca;
  \item versione del codice, configurazione e identificativo dei pesi YOLO
        vengano conservati insieme agli output;
  \item dati vecchi e nuovi non condividano la stessa directory di lavoro.
\end{itemize}

Una consegna minima comprende \filepath{fused_graph.json} e il relativo
\filepath{fused_graph_plot.png}. A seconda del processo possono essere aggiunti
il report di valutazione, \filepath{scene_nodes.json} per la tracciabilità e il
GeoJSON per il passaggio a QGIS. Prima dell'invio esterno, verificare i
metadati: request VGGT, summary e GeoJSON possono contenere percorsi assoluti
della macchina di elaborazione. Cache detection e file \code{.npz} non fanno
normalmente parte della consegna.
